\documentclass[12pt,a4paper]{article}

\usepackage[utf8]{inputenc}
\usepackage[T1]{fontenc}
\usepackage[spanish]{babel}
\usepackage{lmodern}

\usepackage[top=2.5cm, bottom=2.5cm, left=3cm, right=2.5cm]{geometry}

\usepackage[dvipsnames,svgnames,table]{xcolor}
\definecolor{teal600}{HTML}{0F6E56}
\definecolor{teal50} {HTML}{E1F5EE}
\definecolor{purple600}{HTML}{534AB7}
\definecolor{purple50} {HTML}{EEEDFE}
\definecolor{blue600}{HTML}{185FA5}
\definecolor{blue50} {HTML}{E6F1FB}
\definecolor{amber600}{HTML}{BA7517}
\definecolor{amber50} {HTML}{FAEEDA}
\definecolor{coral600}{HTML}{993C1D}
\definecolor{coral50} {HTML}{FAECE7}
\definecolor{gray600}{HTML}{5F5E5A}
\definecolor{gray50} {HTML}{F1EFE8}
\definecolor{codebg} {HTML}{F4F4F4}

\usepackage{tikz}
\usetikzlibrary{
  shapes.geometric,
  arrows.meta,
  positioning,
  fit,
  backgrounds,
  decorations.pathdashing,
  calc
}

\tikzset{

  msvc/.style={
    rectangle, rounded corners=6pt,
    minimum width=3.4cm, minimum height=1.1cm,
    text width=3.2cm, align=center,
    font=\footnotesize\bfseries,
    draw, line width=0.5pt
  },

  msvc-teal/.style={
    msvc,
    fill=teal50, draw=teal600, text=teal600
  },

  msvc-purple/.style={
    msvc,
    fill=purple50, draw=purple600, text=purple600
  },

  msvc-blue/.style={
    msvc,
    fill=blue50, draw=blue600, text=blue600
  },

  msvc-amber/.style={
    msvc,
    fill=amber50, draw=amber600, text=amber600
  },

  msvc-coral/.style={
    msvc,
    fill=coral50, draw=coral600, text=coral600
  },

  msvc-gray/.style={
    msvc,
    fill=gray50, draw=gray600, text=gray600
  },

  actor/.style={
    circle, minimum size=0.7cm,
    fill=teal50, draw=teal600, text=teal600,
    font=\scriptsize\bfseries
  },

  sync/.style={
    -Stealth, thick, draw=blue600
  },

  async/.style={
    -Stealth, thick, dashed, draw=amber600
  },
 
  internal/.style={
    -Stealth, draw=teal600, line width=0.6pt
  },

  container/.style={
    rectangle, rounded corners=8pt,
    dashed, draw=teal600, line width=0.8pt,
    inner sep=10pt
  }
}

\usepackage{booktabs}
\usepackage{tabularx}
\usepackage{array}
\usepackage{multirow}
\renewcommand{\arraystretch}{1.4}

\usepackage{enumitem}
\setlist[itemize]{noitemsep, topsep=4pt}
\setlist[enumerate]{noitemsep, topsep=4pt}

\usepackage{listings}
\lstset{
  backgroundcolor=\color{codebg},
  basicstyle=\ttfamily\footnotesize,
  breaklines=true,
  frame=single,
  rulecolor=\color{gray600},
  numbers=left, numberstyle=\tiny\color{gray600},
  tabsize=2
}


\usepackage[hidelinks, colorlinks=true,
  linkcolor=black, urlcolor=black,
  citecolor=teal600]{hyperref}

\usepackage{parskip}
\usepackage{microtype}
\usepackage{float}
\usepackage{caption}
\captionsetup{font=small, labelfont=bf}

\usepackage{tcolorbox}
\tcbuselibrary{skins,breakable}
\newtcolorbox{infobox}[1][]{
  colback=black50, colframe=black600,
  boxrule=0.5pt, arc=4pt,
  fonttitle=\bfseries\small,
  title={#1}
}
\newtcolorbox{warnbox}[1][]{
  colback=amber50, colframe=amber600,
  boxrule=0.5pt, arc=4pt,
  fonttitle=\bfseries\small,
  title={#1}
}

\begin{document}


\begin{titlepage}
  \centering
  \vspace*{1.5cm}
  {\Large\bfseries\color{black}
    Taller Formativo\\[0.4em]
    Arquitectura de Microservicios\\[0.2em]
    Notación C4 Nivel 3\\[1.5cm]}

  {\large Sistema de Gestión de Solicitudes\\
   de Soporte Técnico de Internet\\[2cm]}

    \begin{flushleft}
        \textbf{Equipo de trabajo:}\\[0.4cm]
        \hspace{1cm}
        \begin{minipage}{0.85\textwidth}
            \begin{itemize}
                \setlength{\itemsep}{0.3cm}

                    \item Cando Moreno Robinson Rodrigo \quad --- \quad Arquitecto Líder            
                \item Luna Mora Ernesto Gregory \quad --- \quad Diseñador de Servicios

                \item Pacheco Cárdenas Cristhian Daniel \quad --- \quad Responsable de Documentación
            \end{itemize}
        \end{minipage}
    \end{flushleft}

    \vspace{1.5cm}

    \begin{flushleft}
        \hspace{1cm}
        \begin{tabular}{ll}
            \textbf{Docente:} & Gleiston Cicerón Guerrero Ulloa \\[0.3cm]
            \textbf{Período Académico:} & Séptimo Semestre \\[0.3cm]
            \textbf{Fecha de Entrega:} & 3 de Junio 2026 \\
        \end{tabular}
    \end{flushleft}

  \vfill
  {\small\color{gray600}\today}
\end{titlepage}


\tableofcontents
\newpage


\section{Descripción del sistema propuesto}

El sistema propuesto es una plataforma web orientada a la gestión estructurada de
solicitudes de soporte técnico relacionadas con problemas en el servicio de Internet,
dentro de entornos organizacionales que requieren trazabilidad, seguridad y eficiencia
en la atención.

\subsection{Perfiles de usuario}

\begin{table}[H]
\centering
\caption{Roles del sistema y sus responsabilidades}
\begin{tabularx}{\textwidth}{lX}
\toprule
\textbf{Rol} & \textbf{Responsabilidades} \\
\midrule
\textbf{Cliente}
  & Registra solicitudes mediante formulario web, adjunta evidencias
    (PNG, JPG, PDF) y realiza seguimiento del estado de sus casos. \\
\textbf{Técnico}
  & Accede únicamente a las solicitudes asignadas; actualiza el estado
    (\emph{Pendiente}, \emph{Resuelta}) y registra acciones con marca
    de fecha, hora y usuario. \\
\textbf{Administrador}
  & Supervisa todas las solicitudes, asigna técnicos manualmente y
    tiene acceso completo a registros de auditoría y seguridad. \\
\textbf{Superusuario}
  & Gestión total de usuarios y roles: crea, modifica y elimina
    cuentas, asigna perfiles. \\
\bottomrule
\end{tabularx}
\end{table}

\section{Arquitectura de microservicios — C4 Nivel 3}

\subsection{Principios de diseño aplicados}

\begin{itemize}
  \item \textbf{SRP} (\emph{Single Responsibility Principle}): cada microservicio
        tiene exactamente una razón para cambiar.
  \item \textbf{Database-per-service}: cada servicio posee su propio almacenamiento;
        ningún servicio accede directamente a la base de datos de otro.
  \item \textbf{API Gateway} como único punto de entrada externo, responsable de
        enrutamiento, autenticación delegada y limitación de tasa.
  \item \textbf{Comunicación síncrona} (REST/gRPC) para operaciones que requieren
        respuesta inmediata; \textbf{comunicación asíncrona} (RabbitMQ) para
        desacoplar productores de consumidores en eventos de dominio.
\end{itemize}

\subsection{Diagrama de componentes — C4 Nivel 3}

\vspace{0.25cm}

\begin{figure}[H]
\centering
\resizebox{\textwidth}{!}{%
\begin{tikzpicture}[node distance=1.4cm and 1.8cm]

  \node[actor] (user) {\textbf{C}};
  \node[below=0.1cm of user, font=\scriptsize\bfseries\color{teal600}]
        {Cliente};

  \node[msvc-purple, right=2.2cm of user] (gw)
        {\textbf{API Gateway}\\[2pt]
         \scriptsize Punto de entrada\\REST/HTTPS};

  \draw[sync] (user) -- node[above,font=\scriptsize\color{blue600}]
        {REST} (gw);

\vspace{0.50cm}

  \begin{scope}[on background layer]
    \node[container,
          fit=(dummy),
          minimum width=10cm, minimum height=7.8cm,
          label={[font=\small\bfseries\color{teal600},
                  inner sep=4pt]north:
                  \guillemotleft microservicio\guillemotright\ Cliente
                  \enspace·\enspace SRP: Gestión de solicitudes del cliente}
         ] (ms-client-box) at (3.2,-3.8) {};
  \end{scope}

  \node[msvc-teal, below=1.8cm of gw] (ctrl)
        {\textbf{Request Controller}\\[2pt]
         \scriptsize Endpoints REST\\validación de payload};

  \node[msvc-teal, right=1.6cm of ctrl] (svc)
        {\textbf{Client Service}\\[2pt]
         \scriptsize Reglas de negocio\\estado solicitudes};

  \node[msvc-teal, below=1.3cm of ctrl] (fval)
        {\textbf{File Validator}\\[2pt]
         \scriptsize PNG/JPG/PDF\\límite de tamaño};

  \node[msvc-teal, below=1.3cm of svc] (evtpub)
        {\textbf{Event Publisher}\\[2pt]
         \scriptsize Publica eventos\\a RabbitMQ};

  \node[msvc-teal, below=1.3cm of fval] (repo)
        {\textbf{Client Repository}\\[2pt]
         \scriptsize Solicitudes\\archivos adjuntos};

  % Flechas internas
  \draw[internal] (gw)    -- node[left, font=\scriptsize\color{teal600}]
        {REST·sync} (ctrl);
  \draw[internal] (ctrl)  -- (svc);
  \draw[internal] (ctrl)  -- (fval);
  \draw[internal] (svc)   -- (evtpub);
  \draw[internal] (fval)  -- (repo);
  \draw[internal] (evtpub)-- (repo);

  \node[msvc-amber, below=1.8cm of evtpub] (rmq)
        {\textbf{RabbitMQ Broker}\\[2pt]
         \scriptsize Cola async\\ticket.created};

  \draw[async] (evtpub) -- node[right, font=\scriptsize\color{amber600}]
        {async} (rmq);

  \node[msvc-gray, right=2.2cm of gw] (auth)
        {\textbf{Auth Service}\\[2pt]
         \scriptsize JWT · roles\\gestión de sesión};

  \draw[sync, draw=gray600] (gw) --
        node[above, font=\scriptsize\color{gray600}]{gRPC·sync}
        (auth);

  \node[msvc-blue, right=2.2cm of svc] (ticket)
        {\textbf{Ticket Service}\\[2pt]
         \scriptsize Crea y asigna\\tickets de soporte};

  \draw[sync] (svc) --
        node[above, font=\scriptsize\color{blue600}]{gRPC·sync}
        (ticket);

  \draw[async] (rmq) |-
        node[near end, right, font=\scriptsize\color{amber600}]{consume}
        (ticket);

  \node[msvc-coral, right=2.2cm of rmq] (notif)
        {\textbf{Notif. Service}\\[2pt]
         \scriptsize Email · push\\SRP};

  \draw[async] (rmq) --
        node[above, font=\scriptsize\color{amber600}]{consume}
        (notif);

  \node[msvc-teal, below=1.4cm of repo] (pg)
        {\textbf{PostgreSQL}\\[2pt]
         \scriptsize DB propia\\del servicio};

  \draw[internal] (repo) -- (pg);

  \node[msvc-teal, left=1.6cm of pg] (minio)
        {\textbf{MinIO / S3}\\[2pt]
         \scriptsize Archivos adjuntos\\PNG/JPG/PDF};

  \draw[internal] (fval) |- (minio);

\end{tikzpicture}
}%
\caption{Diagrama C4 Nivel 3 — Microservicio Cliente del sistema de soporte técnico.}
\label{fig:c4-nivel3}
\end{figure}

\section{Descripción de microservicios}

\subsection{Microservicio Cliente}
\textbf{Responsabilidad única (SRP):}
  Gestionar el ciclo de vida de las solicitudes de soporte enviadas por el usuario
  cliente: registro, validación de archivos, consulta de estado y emisión del evento
  de dominio \texttt{ticket.created}.


\subsubsection{Componentes internos}

\begin{table}[H]
\centering
\caption{Componentes del Microservicio Cliente}
\begin{tabularx}{\textwidth}{lX}
\toprule
\textbf{Componente} & \textbf{Descripción} \\
\midrule
\texttt{Request Controller}
  & Expone los endpoints REST (\texttt{POST /solicitudes},
    \texttt{GET /solicitudes/\{id\}}). Valida el payload
    entrante y delega al servicio de negocio. \\
\texttt{Client Service}
  & Orquesta la lógica de negocio: crea, consulta y actualiza
    solicitudes. Coordina la validación de archivos y la
    publicación de eventos. \\
\texttt{File Validator}
  & Verifica el tipo MIME real (magic bytes), la extensión
    permitida (PNG, JPG, PDF) y el tamaño máximo configurable
    antes de delegar el almacenamiento. \\
\texttt{Event Publisher}
  & Serializa y publica el evento de dominio
    \texttt{ticket.created} en la cola de RabbitMQ.
    Implementa el patrón \emph{outbox transaccional} para
    garantizar entrega al menos una vez. \\
\texttt{Client Repository}
  & Capa de acceso a datos sobre PostgreSQL (propiedad
    exclusiva del servicio). Gestiona solicitudes y metadatos
    de archivos. \\
\bottomrule
\end{tabularx}
\end{table}

\subsection{API Gateway}

Punto de entrada único para todos los clientes externos. Responsable de:
enrutamiento de peticiones, validación de token JWT (delegando al
\texttt{Auth Service} vía gRPC), limitación de tasa (\emph{rate limiting}) y
terminación TLS.

\subsection{Auth Service}

Gestión de autenticación y autorización: emisión y validación de tokens JWT,
control de roles (Cliente, Técnico, Administrador, Superusuario) y gestión de
sesiones. Base de datos propia.

\subsection{Ticket Service}

Recibe el evento \texttt{ticket.created} de RabbitMQ y persiste el ticket en su
propia base de datos. También responde a consultas síncronas del
\texttt{Client Service} mediante gRPC para obtener el estado actual de un ticket.

\subsection{Notif. Service}

Consume el evento \texttt{ticket.created} y envía notificaciones al cliente
(correo electrónico, push) confirmando la recepción de su solicitud.

\section{Canales de comunicación y justificación tecnológica}


\subsection{REST (síncrono)}

\begin{itemize}
  \item \textbf{Uso:} comunicación entre el navegador del cliente y el API Gateway;
        desde el Gateway hacia el \texttt{Request Controller}.
  \item \textbf{Justificación:} estándar ampliamente adoptado, soporte nativo de
        adjuntos \texttt{multipart/form-data}, sencillo de consumir desde cualquier
        frontend web sin dependencias adicionales.
\end{itemize}

\subsection{gRPC (síncrono)}

\begin{itemize}
  \item \textbf{Uso:} API Gateway → Auth Service; Client Service → Ticket Service.
  \item \textbf{Justificación:} baja latencia gracias a Protocol Buffers (binario) y
        HTTP/2. El contrato estricto mediante archivos \texttt{.proto} reduce errores
        de integración entre servicios internos.
\end{itemize}

\subsection{RabbitMQ (asíncrono)}

\begin{itemize}
  \item \textbf{Uso:} Event Publisher → cola \texttt{ticket.created} → Ticket Service,
        Notif. Service.
  \item \textbf{Justificación frente a Kafka:} RabbitMQ es más apropiado para colas
        de trabajo con enrutamiento flexible (exchanges y routing keys), menor latencia
        en mensajes individuales y menor complejidad operacional para este dominio. Kafka
        sería preferible si se requiriera retención de logs y streaming de alto volumen.
  \item \textbf{Garantías:} entrega persistida, acknowledge manual, patrón
        \emph{outbox transaccional} en el publicador.
\end{itemize}
\vspace{0.5cm}
\section{Repositorio en GitHub}

El Código en LaTex de la práctica se encuentra disponible en:
 \url{https://github.com/CristhianP03/Actividades-Grupales-App-Distribuidas.git}

\newpage


\section{Referencias bibliográficas}


\begin{thebibliography}{9}

\bibitem{Taibi2017}
D.~Taibi, V.~Lenarduzzi y C.~Pahl,
``Processes, Motivations, and Issues for Migrating to Microservices
Architectures: An Empirical Investigation,''
\emph{IEEE Cloud Computing}, vol.~4, n.\textordmasculine~5, pp.~22--32, 2017.
DOI: \href{https://doi.org/10.1109/MCC.2017.4250931}
          {10.1109/MCC.2017.4250931}

\bibitem{Dragoni2017}
N.~Dragoni \textit{et al.},
``Microservices: Yesterday, Today, and Tomorrow,''
en \emph{Present and Ulterior Software Engineering},
pp.~195--216, Springer, 2017.
DOI: \href{https://doi.org/10.1007/978-3-319-67425-4_12}
          {10.1007/978-3-319-67425-4\_12}

\bibitem{Cerny2018}
T.~Cerny, M.~J.~Donahoo y M.~Trnka,
``Contextual Understanding of Microservice Architecture,''
\emph{ACM SIGAPP Applied Computing Review}, vol.~17, n.\textordmasculine~4,
pp.~29--45, 2018.
DOI: \href{https://doi.org/10.1145/3183628.3183631}
          {10.1145/3183628.3183631}

\bibitem{Alshuqayran2016}
N.~Alshuqayran, N.~Ali y R.~Evans,
``A Systematic Mapping Study in Microservice Architecture,''
en \emph{Proc. IEEE SOCA 2016}, pp.~44--51, 2016.
DOI: \href{https://doi.org/10.1109/SOCA.2016.15}
          {10.1109/SOCA.2016.15}

\bibitem{Richardson2018}
C.~Richardson,
\emph{Microservices Patterns}.
Manning Publications, 2018.
ISBN: 978-1617294549.

\bibitem{Montesi2016}
F.~Montesi y J.~Weber,
``Circuit Breakers, Discovery, and API Gateways in Microservices,''
\emph{arXiv:1609.05830}, 2016.
Disponible en: \href{https://arxiv.org/abs/1609.05830}
                    {https://arxiv.org/abs/1609.05830}

\end{thebibliography}

\end{document}
